\section{Модели распределения времени остановки}

\subsection{Конечное дискретное распределение}

Базовый вариант --- не навязывать форму распределения и напрямую оценивать вероятности всех возможных значений:
\begin{equation}
  \widehat p_{t,k}(0),\ldots,\widehat p_{t,k}(K),
  \qquad
  \sum_{j=0}^{K}\widehat p_{t,k}(j)=1.
  \label{eq:finite_distribution}
\end{equation}
Параметры получаются из вектора признаков \(G_{t,k}\), построенного по \(Z_{t,k}\), \(Q_{t,k}\) и простым статистикам уверенности:
\begin{equation}
  \widehat p_{t,k}=\operatorname{softmax}(A G_{t,k}+b).
  \label{eq:finite_head}
\end{equation}
Эта модель является основной отправной точкой. Она удобна при небольшом \(K\), не требует предположений о хвосте и сразу дает вероятности прошлого, текущего и будущего оптимума.

\subsection{Сглаженная целевая разметка}

Оракульная метка \(\tau^\star\) может быть слишком жесткой. Если два соседних шага дают почти одинаковое значение \(L_{t,j}+\lambda c(j)\), то принудительный выбор одного шага увеличивает шум обучения. Поэтому вместо точечной метки можно использовать сглаженное целевое распределение
\begin{equation}
  p^\star_{t}(j)=
  \frac{\exp\{-(L_{t,j}+\lambda c(j))/\beta\}}
  {\sum_{r=0}^{K}\exp\{-(L_{t,r}+\lambda c(r))/\beta\}},
  \label{eq:smoothed_target}
\end{equation}
где \(\beta>0\) управляет степенью сглаживания. При малом \(\beta\) модель близка к точечной разметке; при большом \(\beta\) распределение становится мягче.

Целевая функция для конечной модели записывается как перекрестная энтропия:
\begin{equation}
  \mathcal L_{\mathrm{dist}}=-\sum_{t}\sum_{k=0}^{K}\sum_{j=0}^{K}
  p^\star_t(j)\log \widehat p_{t,k}(j).
  \label{eq:distribution_loss}
\end{equation}
Если используется точечная разметка, то \(p^\star_t\) заменяется единичной массой в \(\tau_t^\star\).

\subsection{Смесь геометрических распределений}

Если большинство объектов останавливается рано, а поздние остановки встречаются реже, естественным приближением является смесь геометрических распределений:
\begin{equation}
  \widehat p_{t,k}(j)=\sum_{m=1}^{M}\pi_{t,k,m}(1-\rho_{t,k,m})^j\rho_{t,k,m},
  \qquad j=0,\ldots,K.
  \label{eq:geometric_mixture}
\end{equation}
После ограничения на конечный отрезок \(0,\ldots,K\) распределение нормируется. Компоненты смеси соответствуют разным режимам сложности: простые позиции имеют большую вероятность ранней остановки, сложные --- более медленный спад.

Эта модель использует меньше свободных параметров, чем произвольное дискретное распределение. Она полезна как регуляризованная альтернатива, особенно при малом объеме обучающей выборки для модели времени.

\subsection{Дискретизированная смесь экспоненциальных распределений}

Близкая по смыслу модель получается из непрерывного времени. Для \(j=0,\ldots,K\) можно задать
\begin{equation}
  \widehat p_{t,k}(j)=\sum_{m=1}^{M}\pi_{t,k,m}
  \left[\exp(-\alpha_{t,k,m}j)-\exp(-\alpha_{t,k,m}(j+1))\right],
  \label{eq:exponential_mixture}
\end{equation}
с последующей нормировкой на конечном отрезке. Эта форма удобна, если время остановки трактуется как дискретизация непрерывного времени ожидания.

В работе эта модель рассматривается не как обязательное усложнение, а как один компактный способ проверить, достаточно ли простого экспоненциального хвоста для описания данных.

\subsection{Степенной спад на конечном отрезке}

Отдельное семейство \texttt{power} в реализации соответствует вероятностям с алгебраическим убыванием по индексу шага после нормировки на \(\{0,\ldots,K\}\). Такая форма служит контрастом к экспоненциально-геометрическому спаду: она допускает относительно большую массу на умеренно больших \(j\) и используется как проверка гипотезы о ``мягком'' хвосте без введения второго параметра дисперсии, как в отрицательном биномиальном законе.

\subsection{Отрицательное биномиальное (тяжелый хвост)}

В ряде контекстов поздние остановки могут встречаться чаще, чем предсказывает геометрическая или экспоненциальная модель. Тогда нужна тяжелохвостая форма. В качестве основной такой формы берется отрицательное биномиальное распределение:
\begin{equation}
  \widehat p_{t,k}(j)=
  \binom{j+r_{t,k}-1}{j}(1-\rho_{t,k})^j\rho_{t,k}^{r_{t,k}},
  \qquad j=0,\ldots,K.
  \label{eq:negative_binomial}
\end{equation}
После ограничения на \(0,\ldots,K\) вероятности нормируются. Эта модель гибче геометрической: она допускает большую дисперсию и лучше описывает редкие сложные случаи.

\subsection{Сравнение семейств}

В работе не требуется вводить больше семейств, чем необходимо для проверки основных гипотез. В экспериментах ARC-AGI используются шесть вариантов: конечное дискретное распределение, сглаженная целевая разметка, смесь геометрических распределений, дискретизированная смесь экспоненциальных распределений, степенной спад на отрезке и отрицательное биномиальное распределение. Их роли сведены в табл.~\ref{tab:distribution_families}.

\begin{table}[H]
\centering
\caption{Семейства распределений времени остановки}
\label{tab:distribution_families}
\begin{tabularx}{\textwidth}{lYY}
\toprule
Семейство & Когда полезно & Основной проверяемый эффект \\
\midrule
Конечное дискретное & малое или среднее число шагов & форма распределения без ограничений \\
Сглаженная разметка & несколько близких минимумов риска & устойчивость обучения \\
Смесь геометрических & ранняя остановка преобладает & несколько режимов сложности \\
Смесь экспоненциальных & нужен компактный спад хвоста & связь с временем ожидания \\
Степенной спад (\texttt{power}) & хвост медленнее экспоненты & алгебраическая форма на \(\{0,\ldots,K\}\) \\
Отрицательное биномиальное & заметны поздние остановки & тяжелый хвост, два параметра формы \\
\bottomrule
\end{tabularx}
\end{table}

